% ===========================================================================
%  Projet STORE -- présentation à la direction (public : GRD).
%  Thème GridLab (gridlab-templates/beamer), installé dans TEXMFHOME.
%
%  Chiffres : agrégats tirés des README de lv-screening et injection-scenarios
%  (aucune donnée par réseau ou par nœud). Les \todo{} restent à compléter.
% ===========================================================================
\documentclass[aspectratio=169,10pt]{beamer}
\usetheme{gridlab}

\usepackage[french]{babel}
\usetikzlibrary{babel}   % babel-french rend : ; ! ? actifs, ce qui casse TikZ
\usepackage{colortbl}    % \rowcolor : scénario étudié dans le tableau des scénarios
\usepackage{array}       % >{\raggedright} dans les colonnes p{} des tableaux

\title[Projet STORE]{Projet STORE}
\subtitle{Valeur pour la recherche, collaboration avec OIKEN\\
          et IA dans l'exploitation des réseaux}
\author[F. Sossan]{Prof. Fabrizio Sossan}
\institute{Gridlab, HES-SO Valais-Wallis\\fabrizio.sossan@hevs.ch}
\date{Présentation à la direction -- \today}

% Chiffre clé : \kpi[largeur]{valeur}{légende}
\newcommand{\kpi}[3][2.9cm]{%
  \begin{tikzpicture}
    \node[fill=wash, rounded corners=4pt, text width=#1, minimum height=2.6cm,
          align=center, inner sep=6pt]
      {{\fontsize{24}{26}\selectfont\bfseries\color{primary}#2}\\[6pt]
       {\footnotesize\color{subtle}#3}};
  \end{tikzpicture}%
}

% Partenaire financeur : \partner{nom}{rôle}{en espèces}{en nature}
\newcommand{\partner}[4]{%
  \begin{tikzpicture}
    \node[fill=wash, rounded corners=4pt, text width=3.6cm, minimum height=3.4cm,
          align=center, inner sep=6pt]
      {{\Large\bfseries\color{primary}#1}\\[2pt]
       {\footnotesize\color{subtle}#2}\\[10pt]
       {\fontsize{20}{22}\selectfont\bfseries\color{ink}#3}\\[2pt]
       {\footnotesize\color{subtle}en espèces\\+ #4 en nature}};
  \end{tikzpicture}%
}

% Horizon de temps d'une fonction d'exploitation
\newcommand{\horizon}[1]{{\footnotesize\color{subtle}#1}\par\smallskip}

% Styles des schémas
\tikzset{
  stage/.style={rectangle, rounded corners=3pt, fill=wash, draw=primary,
                line width=0.6pt, align=center, font=\footnotesize},
  stagehi/.style={stage, fill=primary, text=white},
  stageai/.style={stage, fill=secondary, draw=secondary, text=white},
}

\begin{document}

\begin{frame}[plain,noframenumbering]
  \titlepage
\end{frame}

\begin{frame}{Sommaire}
  \tableofcontents
\end{frame}

% ==========================================================================
\section[STORE et la recherche]{Projet STORE : valeur pour la recherche sur les réseaux}
% ==========================================================================

\begin{frame}{Un double objectif}
  \begin{columns}[T,onlytextwidth]
    \begin{column}{0.48\textwidth}
      \begin{block}{1 \quad Objectif principal}
        \textbf{Des outils pour les gestionnaires de réseau.}\par\smallskip
        Les construire ensemble, en assurant le transfert de la recherche
        vers l'exploitation.
      \end{block}
    \end{column}
    \begin{column}{0.48\textwidth}
      \begin{block}{2 \quad Objectif de recherche}
        \textbf{Des outils avancés d'aide à la décision.}\par\smallskip
        Pour la planification et l'exploitation des réseaux de distribution.
      \end{block}
    \end{column}
  \end{columns}

  \vspace{2em}
  \centering
  La recherche alimente les outils ; leur usage sur des réseaux réels
  alimente la recherche.
\end{frame}

\begin{frame}{Économie : un flagship Innosuisse de 6,4\,MCHF}
  \centering
  \kpi{6,4}{MCHF\\coût total du projet}\hfill
  \kpi{4,25}{MCHF d'Innosuisse\\overhead compris}\hfill
  \kpi{2,7}{MCHF apportés par\\les partenaires}\hfill
  \kpi{3,05}{MCHF à la HEI-VS\\hors overhead}

  \vspace{1.4em}
  \raggedright
  Décembre 2023 -- juillet 2027 (44 mois). Quatre hautes écoles (ETHZ, EPFL,
  BFH, HEI-VS) et treize partenaires de mise en œuvre ; la HEI-VS coordonne.

  \vspace{0.8em}
  {\footnotesize\color{subtle} Montants de la demande Innosuisse 108.230
  FS-EE. Partenaires : 0,97\,MCHF en espèces et 1,76\,MCHF en nature. HEI-VS :
  2,08\,MCHF couverts par Innosuisse et 0,97\,MCHF versés par les
  partenaires.}
\end{frame}

\begin{frame}{Économie : les gestionnaires de réseau cofinancent}
  \centering
  \partner{OIKEN}{distribution\\Sion et Sierre}{246\,kCHF}{294\,kCHF}\hfill
  \partner{Valgrid}{réseau suprarégional\\Valais}{90\,kCHF}{24\,kCHF}\hfill
  \partner{Swissgrid}{réseau de transport\\Suisse}{200\,kCHF}{50\,kCHF}

  \vspace{1.4em}
  \raggedright
  Ensemble, 536\,kCHF en espèces : 55\,\% des espèces apportées par
  l'ensemble des partenaires. Swissgrid et Valgrid financent la
  planification coordonnée du stockage, du transport à la distribution, sur
  le cas du Valais central avec OIKEN.

  \vspace{0.8em}
  {\footnotesize\color{subtle} Montants de la demande Innosuisse 108.230
  FS-EE ; en nature : heures de travail des partenaires.}
\end{frame}

\begin{frame}{Développer pleinement les compétences locales}
  \centering
  \begin{tikzpicture}[node distance=0.7cm,
      every node/.style={text width=3.4cm, minimum height=1.5cm}]
    \node (hire)  [stage]                 {Recruter des ingénieurs au profil
                                           d'exploitant};
    \node (train) [stage, right=of hire]  {Les former sur les réseaux d'OIKEN
                                           et au GridLab};
    \node (keep)  [stagehi, right=of train]{Des compétences qui restent en
                                            Valais};
    \draw [arrow, primary] (hire)  -- (train);
    \draw [arrow, primary] (train) -- (keep);
  \end{tikzpicture}

  \vspace{1.8em}
  \raggedright
  Ces recrues renforcent aussi le GridLab : essais, démonstrateurs et
  transfert vers les gestionnaires de réseau.
\end{frame}

% ==========================================================================
\section[Collaboration OIKEN]{Collaboration avec OIKEN : résultats marquants}
% ==========================================================================

\begin{frame}[c]
  \centering
  {\small\color{subtle}La question}\par
  \vspace{0.8em}
  {\LARGE\bfseries\color{primary}Les réseaux d'OIKEN sont-ils prêts\\
   pour la transition énergétique ?}\par
  \vspace{0.8em}
  {\color{secondary}\rule{0.3\textwidth}{1.3pt}}\par
  \vspace{0.8em}
  Une question ambitieuse : nous y répondons réseau par réseau, à partir des
  données réelles.
\end{frame}

\begin{frame}{Une analyse bottom-up, unique}
  \begin{columns}[c,onlytextwidth]
    \begin{column}{0.47\textwidth}
      \centering
      \begin{tikzpicture}[node distance=0.55cm,
          every node/.style={text width=5cm, minimum height=1.1cm}]
        \node (node) [stage] {\textbf{Données réseau d'OIKEN}\\chaque
                              raccordement : demande, PV, VE};
        \node (lv)   [stage, above=of node] {\textbf{1116 réseaux BT}\\calcul
                              de répartition de charge, pas de 15 min};
        \node (hv)   [stagehi, above=of lv] {\textbf{11 postes HT/MT}\\somme
                              des réseaux BT alimentés};
        \draw [arrow, primary] (node) -- (lv);
        \draw [arrow, primary] (lv)   -- (hv);
      \end{tikzpicture}
    \end{column}
    \begin{column}{0.49\textwidth}
      \begin{itemize}
        \item \textbf{Unique} : nous partons des données réelles du réseau,
              raccordement par raccordement, et non de réseaux types
        \item \textbf{Scénarios} : scénario-cadre de l'OFEN pour la
              planification des réseaux, horizons 2030 et 2040
        \item \textbf{Traçable} : chaque dépassement est rattaché à son
              réseau, à son nœud et à sa cause (demande, PV, VE)
        \item \textbf{Rapide} : tout le portefeuille en \(\approx 30\) s
              sur un ordinateur portable
      \end{itemize}
    \end{column}
  \end{columns}
\end{frame}

\begin{frame}{Le défi : un grand jeu de données}
  \centering
  \kpi[3.6cm]{1116}{réseaux basse tension,\\avec le réseau MT}\hfill
  \kpi[3.6cm]{3,9\,GB}{base de données réseau\\extraite et reconstruite}\hfill
  \kpi[3.6cm]{3300\,km}{de lignes BT,\\100\,000 tronçons}

  \vspace{1.4em}
  \raggedright
  Extraire les réseaux de la base du GRD, reconstruire la topologie BT et
  MT, puis contrôler automatiquement chaque réseau avant tout calcul :
  1036 réseaux sur 1116 sont calculables d'emblée, les autres sont signalés
  avec leur cause.
\end{frame}

\begin{frame}{Ce que les contrôles ont révélé}
  \centering
  \footnotesize
  \renewcommand{\arraystretch}{1.3}
  \begin{tabular}{@{}>{\raggedright\arraybackslash}p{8.4cm}r
                     >{\raggedright\arraybackslash}p{3.6cm}@{}}
    \toprule
    Constat & Réseaux & Détail \\
    \midrule
    Réseaux BT maillés (boucles ou branches parallèles), alors que la BT
      est exploitée en radial                                 & 259 & 23\,\% du portefeuille \\
    Section de conducteur manquante ou fictive (0 ou 999\,mm²) & 592 & 2721 tronçons, 68\,km \\
    \og PV existant \fg{} : 12 nœuds BT portent 183\,MW, chacun au-delà de
      la puissance de son propre transformateur, ce qu'aucun vrai
      raccordement BT ne permet (centrales, pompages)          & 11  & 183\,MW sur 299\,MW \\
    Réseaux non calculables : sans transformateur, nœuds isolés, ligne
      bouclée sur un seul nœud                                 & 9   & 2 + 2 + 5 réseaux \\
    Transformateurs fictifs (10\,kVA ou moins) ou saisis à l'envers
      (BT vers MT)                                             & 36  & 13 fictifs, 30 à l'envers \\
    Rattachement au poste HT/MT absent                         & 18  & zone MT erronée pour
                                                                       environ 30\,\% des réseaux \\
    Réseaux sans aucun raccordement                            & 35  & \\
    \bottomrule
  \end{tabular}

  \vspace{1em}
  \raggedright
  \small
  Valeur pour OIKEN : une liste de corrections concrètes, réseau par réseau,
  pour fiabiliser la base avant toute étude de planification.
\end{frame}

\begin{frame}{Les scénarios : le scénario-cadre de la Confédération}
  Le \textbf{scénario-cadre pour la planification des réseaux} (OFEN, issu
  des Perspectives énergétiques 2050+) est la base commune de la
  Confédération pour planifier les réseaux électriques. Il retient trois
  scénarios :

  \vspace{1em}
  \centering
  \small
  \renewcommand{\arraystretch}{1.35}
  \begin{tabular}{@{}lcccc@{}}
    \toprule
    & \multicolumn{2}{c}{PV installé en Suisse (GW)}
    & \multicolumn{2}{c}{VE en Suisse (millions)} \\
    \cmidrule(lr){2-3}\cmidrule(l){4-5}
    Scénario & 2030 & 2040 & 2030 & 2040 \\
    \midrule
    \rowcolor{wash}
    \textbf{Référence} (ZERO Base) & \textbf{9,8} & \textbf{24,1}
                                   & \textbf{0,93} & \textbf{2,94} \\
    Divergence (ZERO A)            & 7,7  & 10,1 & 0,98 & 3,23 \\
    Couplage des secteurs (ZERO B) & 12,2 & 30,1 & 0,87 & 2,52 \\
    \bottomrule
  \end{tabular}

  \vspace{1.2em}
  \raggedright
  Résultats présentés : \textbf{scénario Référence, horizon 2040}, ramené à
  Sion et Sierre au prorata de la population, soit environ 266\,MW de PV et
  32\,500 VE. Divergence et Couplage des secteurs : prochaine étape, la
  chaîne est prête.
\end{frame}

\begin{frame}[fragile]{Premiers éléments de réponse : presque prêts en 2030, pas en 2040}
  \begin{columns}[c,onlytextwidth]
    \begin{column}{0.52\textwidth}
      \centering
      {\small\textbf{Réseaux BT en dépassement, par type}\\
       \footnotesize\color{subtle}ordres de grandeur}\par
      \vspace{0.4em}
      \shorthandoff{;:!?}% pgfplots lit le ; final avant que babel le rende
      \begin{tikzpicture}
        \begin{axis}[
            xbar, scale only axis, width=2.9cm, height=3.6cm,
            bar width=6pt, enlarge y limits=0.18,
            xmin=0, xmax=300, clip=false, hide x axis,
            axis y line*=left, y axis line style={draw=subtle},
            ytick=data, symbolic y coords={over,under,trafo,line},
            yticklabels={Surtension, Sous-tension, Transformateur surchargé,
                         Ligne surchargée},
            yticklabel style={font=\footnotesize, text=ink},
            major tick length=0pt,
            nodes near coords={\pgfmathprintnumber[assume math mode=true]{\pgfplotspointmeta}},
            nodes near coords style={font=\scriptsize, text=ink},
            area legend, reverse legend,
            legend style={draw=none, font=\footnotesize, text=ink,
                          at={(1,1.03)}, anchor=south east, legend columns=2,
                          /tikz/every even column/.append style={column sep=8pt}},
          ]
          \addplot[fill=primary, draw=none]
            coordinates {(30,over) (216,under) (241,trafo) (270,line)};
          \addplot[fill=primary!40, draw=none]
            coordinates {(0,over) (30,under) (3,trafo) (28,line)};
          \legend{2040, 2030}
        \end{axis}
      \end{tikzpicture}
    \end{column}
    \begin{column}{0.44\textwidth}
      \kpi[2.2cm]{4\,\%}{des réseaux BT\\en 2030\\(45 sur 1036)}\hfill
      \kpi[2.2cm]{36\,\%}{des réseaux BT\\en 2040\\(371 sur 1036)}

      \vspace{0.8em}
      En 2030, 12 de ces réseaux dépassent déjà avec la demande actuelle ; la
      recharge des VE en ajoute 33, le PV aucun.

      \vspace{0.4em}
      De 2030 à 2040 : 3\,\texttimes{} plus de VE, 2,4\,\texttimes{} plus de
      PV, 8\,\texttimes{} plus de réseaux en dépassement.
    \end{column}
  \end{columns}

  \vspace{0.8em}
  {\footnotesize\color{subtle} Scénario-cadre, scénario Référence (ZERO Base)
  2030 et 2040, ramené à Sion et Sierre au prorata de la population ; journée
  représentative au pas de 15 min. Ordres
  de grandeur : les comptes varient avec le tirage aléatoire de la demande et
  des arrivées de VE.}
\end{frame}

\begin{frame}[fragile]{D'où viennent les dépassements ? De la demande, bien plus que du PV}
  \begin{columns}[c,onlytextwidth]
    \begin{column}{0.52\textwidth}
      \centering
      {\small\textbf{Réseaux BT en dépassement en 2040, par cause}}\par
      \vspace{0.4em}
      \shorthandoff{;:!?}% pgfplots lit le ; final avant que babel le rende
      \begin{tikzpicture}
        \begin{axis}[
            xbar, scale only axis, width=2.9cm, height=3.6cm,
            bar width=6pt, enlarge y limits=0.18,
            xmin=0, xmax=300, clip=false, hide x axis,
            axis y line*=left, y axis line style={draw=subtle},
            ytick=data, symbolic y coords={over,under,trafo,line},
            yticklabels={Surtension, Sous-tension, Transformateur surchargé,
                         Ligne surchargée},
            yticklabel style={font=\footnotesize, text=ink},
            major tick length=0pt,
            nodes near coords={\pgfmathprintnumber[assume math mode=true]{\pgfplotspointmeta}},
            nodes near coords style={font=\scriptsize, text=ink},
            area legend, reverse legend,
            legend style={draw=none, font=\footnotesize, text=ink,
                          at={(1,1.03)}, anchor=south east, legend columns=2,
                          /tikz/every even column/.append style={column sep=8pt}},
          ]
          \addplot[fill=secondary, draw=none]
            coordinates {(30,over) (0,under) (8,trafo) (38,line)};
          \addplot[fill=primary, draw=none]
            coordinates {(0,over) (216,under) (241,trafo) (268,line)};
          \legend{PV, Demande (ménages et VE)}
        \end{axis}
      \end{tikzpicture}
    \end{column}
    \begin{column}{0.44\textwidth}
      \kpi[2.2cm]{357}{réseaux :\\recharge des VE,\\le soir}\hfill
      \kpi[2.2cm]{50}{réseaux :\\PV, à midi}

      \vspace{0.8em}
      Demande actuelle des ménages seule : 12 réseaux. Sur les 50 réseaux
      touchés par le PV, 48 le sont aussi par les VE.

      \vspace{0.4em}
      Levier principal pour le GRD : la recharge pilotée des VE.
    \end{column}
  \end{columns}

  \vspace{0.8em}
  {\footnotesize\color{subtle} Load flow avec la demande seule, puis
  demande + PV et demande + VE : une cause est retenue si, ajoutée à la
  demande, elle suffit à créer le dépassement. Un réseau peut avoir les deux
  causes (PV à midi, VE le soir) ; aucun dépassement n'exige les deux à la fois.}
\end{frame}

\begin{frame}{L'effet remonte jusqu'aux postes HT/MT}
  \centering
  \kpi[4.6cm]{0,6 à 2,0\,\texttimes}{pointe de recharge des VE le soir,\\
    rapportée à la puissance des transformateurs HT/MT}\hspace{1.5cm}%
  \kpi[4.6cm]{60\,\%}{export PV à midi (environ),\\
    rapporté à cette même puissance}

  \vspace{1.6em}
  \raggedright
  Pour le GRD, la planification BT seule ne suffit pas : il faut coordonner
  BT et MT et mobiliser la flexibilité (recharge pilotée, stockage).

  \vspace{0.8em}
  {\footnotesize\color{subtle} Scénario Référence 2040, somme des injections
  des réseaux BT par poste HT/MT. Non modélisés : lignes MT et leurs pertes, clients et
  producteurs raccordés en MT.}
\end{frame}

% ==========================================================================
\section[IA et exploitation]{L'IA dans la boucle de l'exploitation des réseaux}
% ==========================================================================

\begin{frame}{Un grand potentiel, sur les trois fonctions d'exploitation}
  \begin{columns}[T,onlytextwidth]
    \begin{column}{0.31\textwidth}
      \begin{block}{Planification}
        \horizon{années à mois}
        Explorer des scénarios, fiabiliser les données, préparer les
        décisions d'investissement.
      \end{block}
    \end{column}
    \begin{column}{0.31\textwidth}
      \begin{block}{Programmation}
        \horizon{veille et infrajournalier}
        Prévoir la demande et le PV, programmer le stockage et la recharge
        des VE.
      \end{block}
    \end{column}
    \begin{column}{0.31\textwidth}
      \begin{block}{Contrôle}
        \horizon{temps réel}
        Proposer des consignes, vérifiées par un modèle physique avant
        d'être appliquées.
      \end{block}
    \end{column}
  \end{columns}

  \vspace{2em}
  \centering
  {\large\bfseries\color{primary}L'IA propose, la physique vérifie,
   l'opérateur décide.}
\end{frame}

\begin{frame}{Au GridLab : une GeForce dans la boucle}
  \begin{columns}[c,onlytextwidth]
    \begin{column}{0.5\textwidth}
      \centering
      \begin{tikzpicture}[every node/.style={text width=2.6cm,
                                             minimum height=1.3cm}]
        \node (meas) [stage]   at (0,0)     {Mesures du laboratoire};
        \node (ai)   [stageai] at (3.6,0)   {IA sur GPU GeForce, sur site};
        \node (phys) [stage]   at (3.6,-2.4){Vérification physique};
        \node (act)  [stagehi] at (0,-2.4)  {Consignes aux équipements};
        \draw [arrow, primary] (meas) -- (ai);
        \draw [arrow, primary] (ai)   -- (phys);
        \draw [arrow, primary] (phys) -- (act);
        \draw [arrow, primary] (act)  -- (meas);
      \end{tikzpicture}
    \end{column}
    \begin{column}{0.46\textwidth}
      \begin{itemize}
        \item \textbf{Matériel grand public} : une carte graphique GeForce,
              un coût d'entrée modeste
        \item \textbf{Sur site} : le modèle tourne localement, les données ne
              quittent pas le laboratoire
        \item \textbf{En boucle fermée} : \todo{tâche de l'IA, équipements
              pilotés, résultats}
      \end{itemize}
    \end{column}
  \end{columns}

  \vspace{1.2em}
  Le même schéma, transposé chez un GRD, garde ses données chez lui.
\end{frame}

\def\closingnote{Questions ?}
\closingslide[Merci de votre attention]

\end{document}
